\chapter{Literature Review}\label{chap:background}
Five recent articles (2020--2025) on UAV--UGV task allocation and routing are reviewed below.
 
\textbf{Ma et al. (2024) \cite{ma}, \emph{Drones}.} \textbf{Formulation.} MILP jointly assigning targets and planning trajectories for a UAV/UGV detection team; binary variables mark which agent detects which target and when, continuous variables give position, velocity and acceleration. \textbf{Objective.} Weighted mission completion time, total operating time, energy and trajectory smoothness. \textbf{Constraints.} UAV/UGV kinematics, target assignment, collision avoidance, obstacles/prohibited regions. \textbf{Algorithm.} IGA-B\&B: a GA assigns and sequences targets, A* plans paths, and the GA solution seeds Branch-and-Bound as its initial upper bound. \textbf{Results/metrics.} About 30\% less computation time than standard B\&B; metrics: objective value, computation time; an outdoor test confirmed feasibility. \textbf{Accomplishments.} Unified allocation and path planning. \textbf{Shortcomings.} Simplifying assumptions; cost grows with fleet size and targets.
 
\textbf{Mondal et al. (2025) \cite{mondal}, \emph{J.\ Intell.\ Robot.\ Syst.}} \textbf{Formulation.} Fuel-constrained UAVs serve surveillance tasks with slower UGVs as mobile refueling stations (E-VRPTW); decisions: UGV refueling/rendezvous points and route, UAV task allocation and visit sequences. \textbf{Objective.} Minimize routing cost and mission time while completing all tasks. \textbf{Constraints.} Task coverage, UAV fuel capacity, rendezvous timing, time windows, UAV/UGV speed differences. \textbf{Algorithm.} Sequential: minimum set-cover picks UGV stops, the UGV route is planned, then UAV allocation and E-VRPTW are solved with greedy and constraint-programming heuristics. \textbf{Results/metrics.} 30 scenarios on three scales; fuel-capacity effect, greedy vs.\ constraint programming, UGV-only baseline; metrics: routing cost, mission time, computation time. \textbf{Accomplishments.} Integrates allocation, routing and fuel limits at scale. \textbf{Shortcomings.} UGV-first decomposition loses global optimality; heuristic; computational validation only.
 
\textbf{Chen et al. (2025) \cite{chen}, \emph{Unmanned Systems}.} \textbf{Formulation.} Medical rescue combining allocation, trajectory generation and sequencing; decisions: vehicle per task, collision-free paths, visiting order. \textbf{Objective.} Reduce completion time, travel distance and energy; balance workload. \textbf{Constraints.} Capability limits, task allocation, obstacle avoidance, dynamic feasibility, base start/end, smoothness. \textbf{Algorithm.} ``Allocate--Plan--Optimize'': Enhanced GA, Informed-RRT*, CMA-ES, waypoint post-processing. \textbf{Results/metrics.} 15 tasks: 26.7 min, over 73\% better than K-Means and random; CMA-ES cut distance by 15.1\%; metrics: time, distance, trajectory efficiency, robustness, scalability. \textbf{Accomplishments.} Integrates allocation, planning and sequencing. \textbf{Shortcomings.} No global optimum; simulation only; no dynamic tasks or communication limits.
 
\textbf{Gao et al. (2020) \cite{gao}, \emph{IEEE Access}.} \textbf{Formulation.} Nested VRP (MILP): a UGV carrying a UAV delivers emergency resources; decisions: UGV route, UAV sorties, customer split. \textbf{Objective.} Minimize total service time. \textbf{Constraints.} Each customer served once by UAV or UGV; UAV payload, range, battery; launch/retrieval synchronization and times. \textbf{Algorithm.} ``UGV first, UAV second'': Lin--Kernighan route, Clarke--Wright savings, saving-based iterative improvement. \textbf{Results/metrics.} 60 instances (10/50/100 customers, 20 each); gains of 16.3\%, 14.4\%, 10.4\%; metrics: service time, improvement \%. \textbf{Accomplishments.} Larger instances without exact optimization. \textbf{Shortcomings.} UGV-dominant route limits optimality; single UGV/UAV only (multiple vehicles, time windows, no-fly zones left as future work).
 
\textbf{Wang et al. (2024) \cite{wang}, \emph{Swarm Evol.\ Comput.}} \textbf{Formulation.} Bi-objective allocation and routing of $m$ weeding robots and $l$ spraying drones over $n$ locations (weeding, then spraying) from a depot; binary $xr_{ij},xd_{ij}=1$ if a robot/drone visits $j$ right after $i$ (0 = depot). \textbf{Objectives.} $\min C=\max_k T_k$ ($T_k$: return time of vehicle $k$) and $\min R=\sum_{i=1}^{n}DQ_i\,xd_{i0}$ ($DQ_i$: herbicide left after spraying at $i$). \textbf{Constraints.} $\sum_{i=0}^{n}xr_{ij}=\sum_{i=0}^{n}xd_{ij}=1\ \forall j$ (matching outgoing constraints); depot start/end; subtour elimination; precedence $STD_j=\max(ATD_j,DTR_j)$; refill if $DQ_i<q_j$; balance $DQ_j=AQ_j-q_j$. \textbf{Algorithm.} MOTLBO: priority-based initialization, multi-teacher learning, order crossover, multi-neighborhood search, non-dominated sorting, Pareto archive. \textbf{Results/metrics.} 150 calibration and 300 test instances (40--90 tasks, 3--7 vehicles per type), five runs, 4 s limit; beat six algorithms with GD $=0.0831$, hypervolume $=0.9898$, Spread $=0.6552$; metrics: GD, hypervolume, Spread, C-metric. \textbf{Accomplishments.} Integrates allocation, routing, precedence, replenishment, load-dependent drone speed. \textbf{Shortcomings.} Simplified static obstacles; no battery constraint; approximate reference front.
 
\textbf{Synthesis.} Hybrid, decomposition-based meta-heuristics (GA variants, CMA-ES, TLBO, savings heuristics) on MILP/VRP models dominate. Benchmarks for comparison: 30 scenarios \cite{mondal}, 60 instances \cite{gao}, the 15-task case \cite{chen}, 450 instances \cite{wang}. Metrics: mission time, distance, energy/fuel, computation time; GD, hypervolume and Spread for multi-objective cases. Gaps: sequential decomposition without global optimality, simulation-only validation, limited energy, dynamic and large-fleet treatment.

%\let\cleardoublepage\clearpage